% Part: methods
% Chapter: induction
% Section: inductive-definitions

\documentclass[../../../include/open-logic-section]{subfiles}

\begin{document}

\olfileid[id]{mth}{ind}{idf}

\olsection{Definisi Induktif}

Dalam logika, kita sangat sering mendefinisikan jenis-jenis objek
\emph{secara induktif}, yakni dengan menentukan kaidah-kaidah mengenai apa
yang terhitung sebagai objek dari jenis yang hendak didefinisikan; kaidah
itu menjelaskan cara memperoleh objek baru dari jenis tersebut berdasarkan
objek lama dari jenis yang sama. Misalnya, kita sering mendefinisikan jenis
khusus barisan simbol, seperti suku dan !!{formula} suatu bahasa, melalui
induksi. Sebagai contoh sederhana, perhatikan untai yang terdiri atas huruf
$\mathrm{a}$, $\mathrm{b}$, $\mathrm{c}$, $\mathrm{d}$, simbol~$\circ$,
serta kurung siku $[$ dan $]$, misalnya
``$[[\mathrm{c} \circ \mathrm{d}][$'', ``$[\mathrm{a}[]\circ]$'',
``$\mathrm{a}$'', atau ``$[[\mathrm{a} \circ \mathrm{b}]\circ
\mathrm{d}]$''. Barangkali Anda merasa bahwa ada sesuatu yang ``salah''
pada dua untai pertama: kurung siku sama sekali tidak ``seimbang'' pada
untai pertama, dan Anda mungkin merasa bahwa ``$\circ$'' seharusnya
``menghubungkan'' ungkapan-ungkapan yang masing-masing masuk akal. Untai
ketiga dan keempat tampak lebih baik: bagi setiap ``$[$'' terdapat
``$]$'' penutup (jika memang ada kurung sama sekali), dan bagi setiap
$\circ$ kita dapat menemukan ungkapan yang ``rapi'' pada kedua sisinya,
dikelilingi sepasang kurung siku.

Kita ingin menentukan secara tepat apa yang terhitung sebagai ``suku
rapi''. Pertama-tama, setiap huruf dengan sendirinya merupakan suku rapi.
Apa pun yang bukan sekadar satu huruf harus berbentuk ``$[t \circ s]$'',
dengan $s$ dan $t$ sendiri merupakan suku rapi. Sebaliknya, jika $t$ dan
$s$ rapi, kita dapat membentuk suku rapi baru dengan menempatkan $\circ$ di
antara keduanya dan mengelilinginya dengan sepasang kurung siku. Kita dapat
menggunakan operasi-operasi ini untuk \emph{mendefinisikan} himpunan suku
rapi. Inilah suatu \emph{definisi induktif}.

\begin{defn}[Suku rapi]
  Himpunan \emph{suku rapi} didefinisikan secara induktif sebagai berikut:
  \begin{enumerate}
  \item Setiap huruf $\mathrm{a}$, $\mathrm{b}$, $\mathrm{c}$,
    $\mathrm{d}$ merupakan suku rapi.
  \item Jika $s_1$ dan $s_2$ merupakan suku rapi, maka
    $[s_1 \circ s_2]$ juga merupakan suku rapi.
  \item Tidak ada yang lain yang merupakan suku rapi.
  \end{enumerate}
\end{defn}

Definisi ini memberi tahu kita bahwa sesuatu terhitung sebagai suku rapi
jika dan hanya jika sesuatu itu dapat dibangun menurut dua syarat (1)
dan~(2) dalam sejumlah berhingga langkah. Pada langkah pertama, kita
membangun semua suku rapi yang hanya terdiri atas satu huruf, yakni
\[
\mathrm{a}, \mathrm{b}, \mathrm{c}, \mathrm{d}
\]
Pada langkah kedua, kita menerapkan (2) pada suku-suku yang telah dibangun.
Kita memperoleh
\[
[\mathrm{a} \circ \mathrm{a}], [\mathrm{a} \circ \mathrm{b}],
[\mathrm{b} \circ \mathrm{a}], \dots, [\mathrm{d} \circ \mathrm{d}]
\]
untuk semua kombinasi dua huruf. Pada langkah ketiga, kita menerapkan (2)
lagi pada dua suku rapi mana pun yang sudah dibangun. Kita memperoleh suku
rapi baru seperti $[\mathrm{a} \circ [\mathrm{a} \circ
    \mathrm{a}]]$---di sini $t$ adalah $\mathrm{a}$ dari langkah~1 dan $s$
adalah $[\mathrm{a} \circ \mathrm{a}]$ dari langkah~2---serta
$[[\mathrm{b} \circ \mathrm{c}] \circ [\mathrm{d} \circ \mathrm{b}]]$,
yang dibangun dari dua suku $[\mathrm{b} \circ \mathrm{c}]$ dan
$[\mathrm{d} \circ \mathrm{b}]$ dari langkah~2. Demikian seterusnya.
Klausa (3) mencegah apa pun yang tidak dibangun dengan cara ini menyusup ke
dalam himpunan suku rapi.

Perhatikan bahwa kita belum membuktikan bahwa setiap barisan simbol yang
``terasa'' rapi memang rapi menurut definisi ini. Namun, semestinya jelas
bahwa semua yang dapat kita bangun memang ``terasa rapi'': kurung sikunya
seimbang, dan $\circ$ menghubungkan bagian-bagian yang dengan sendirinya
rapi.

Ciri utama definisi induktif ialah bahwa, jika Anda hendak membuktikan
sesuatu mengenai semua suku rapi, definisi tersebut memberi tahu Anda kasus
mana saja yang harus dipertimbangkan. Misalnya, jika diberi tahu bahwa $t$
adalah suku rapi, definisi induktif memberi tahu Anda seperti apa bentuk
$t$: $t$ dapat berupa sebuah huruf, atau dapat berupa $[s_1 \circ s_2]$
untuk sepasang suku rapi $s_1$ dan~$s_2$. Karena klausa (3), hanya itulah
kemungkinannya.

Ketika membuktikan klaim mengenai seluruh anggota suatu himpunan yang
didefinisikan secara induktif, bentuk induksi kuat menjadi sangat penting.
Misalnya, andaikan kita hendak membuktikan bahwa bagi setiap suku rapi
sepanjang~$n$, banyaknya $[$ di dalamnya adalah~$< n/2$. Ini dapat
dipandang sebagai klaim mengenai semua~$n$: bagi setiap $n$, banyaknya $[$
dalam sebarang suku rapi sepanjang~$n$ adalah $< n/2$.

\begin{prop}
  Untuk sebarang $n$, banyaknya $[$ dalam suatu suku rapi sepanjang~$n$
  adalah $< n/2$.
\end{prop}

\begin{proof}
Untuk membuktikan hasil ini melalui induksi (kuat), kita harus menunjukkan
bahwa klaim kondisional berikut benar:
\begin{quote}
  Jika bagi setiap $l < k$, sebarang suku rapi sepanjang~$l$ memiliki
  $< l/2$ buah $[$, maka sebarang suku rapi sepanjang~$k$ memiliki
  $< k/2$ buah $[$.
\end{quote}
Untuk menunjukkan kondisional ini, andaikan antesedennya benar, yakni
andaikan bahwa bagi sebarang $l<k$, suku rapi sepanjang~$l$ memuat
$< l/2$ buah $[$. Asumsi ini kita sebut hipotesis induksi. Kita hendak
menunjukkan bahwa hal yang sama benar bagi suku rapi sepanjang~$k$.

Jadi, andaikan $t$ adalah suku rapi sepanjang~$k$. Karena suku rapi
didefinisikan secara induktif, kita mempunyai dua kasus: (1)~$t$ hanya
berupa satu huruf, atau (2)~$t$ adalah $[s_1 \circ s_2]$ untuk suatu suku
rapi $s_1$ dan~$s_2$.
\begin{enumerate}
\item $t$ adalah sebuah huruf. Maka $k = 1$, dan banyaknya $[$ dalam $t$
adalah~$0$. Karena $0 < 1/2$, klaim tersebut berlaku.
\item $t$ adalah $[s_1 \circ s_2]$ untuk suatu suku rapi $s_1$ dan~$s_2$.
  Misalkan $l_1$ adalah panjang~$s_1$ dan $l_2$ adalah panjang~$s_2$.
  Maka panjang~$k$ dari $t$ adalah $l_1+l_2+3$ (panjang $s_1$ dan $s_2$
  ditambah tiga simbol $[$, $\circ$, $]$). Karena $l_1+l_2+3$ selalu lebih
  besar daripada $l_1$, berlaku $l_1 < k$. Demikian pula, $l_2 < k$. Ini
  berarti hipotesis induksi berlaku pada suku $s_1$ dan~$s_2$:
  bilangan~$m_1$ yang menyatakan banyaknya $[$ dalam $s_1$ adalah
  $< l_1/2$, dan
  bilangan~$m_2$ yang menyatakan banyaknya $[$ dalam~$s_2$ adalah
  $< l_2/2$.

  Banyaknya $[$ dalam $t$ adalah banyaknya $[$ dalam~$s_1$, ditambah
  banyaknya $[$ dalam~$s_2$, ditambah~$1$, yakni $m_1 + m_2 + 1$. Karena
  $m_1 < l_1/2$ dan $m_2 < l_2/2$, kita memperoleh:
  \[
  m_1 + m_2 + 1 < \frac{l_1}{2} + \frac{l_2}{2} + 1 = \frac{l_1+l_2+2}{2} < \frac{l_1+l_2+3}{2} = k/2.
  \]
\end{enumerate}
Dalam setiap kasus, kita telah menunjukkan bahwa banyaknya $[$ dalam $t$
adalah $< k/2$ (berdasarkan hipotesis induksi). Melalui induksi kuat,
proposisi tersebut mengikuti.
\end{proof}

\begin{prob}
  Definisikan himpunan suku superrapi dengan ketentuan berikut:
  \begin{enumerate}
  \item Setiap huruf $\mathrm{a}$, $\mathrm{b}$, $\mathrm{c}$,
    $\mathrm{d}$ merupakan suku superrapi.
  \item Jika $s$ merupakan suku superrapi, maka $[s]$ juga merupakan suku
    superrapi.
  \item Jika $s_1$ dan $s_2$ merupakan suku superrapi, maka
    $[s_1 \circ s_2]$ juga merupakan suku superrapi.
  \item Tidak ada yang lain yang merupakan suku superrapi.
  \end{enumerate}
  Tunjukkan bahwa banyaknya $[$ dalam suatu suku superrapi~$t$ sepanjang
  $n$ adalah $\le n/2 +1$.
\end{prob}

\end{document}
